\chapter{Frontend}
\label{cap:frontend}

\begin{edital}
Tecnologias e práticas frontend web --- HTML, CSS, UX, Ajax, frameworks
(VueJS, Angular, React); padrões de frontend; SPA e PWA; UX (sistemas de
gestão de conteúdo, arquitetura de informação, portais, workflow,
acessibilidade, usabilidade).
\end{edital}
\peso{MÉDIO --- metade é distinguir tecnologias (SPA $\times$ PWA,
frameworks, mobile) e metade é \textbf{ler código HTML/CSS/JS} e dizer a
saída. Não dá para decorar só conceito; treine ler snippet.}

\section{HTML e CSS}

\textbf{HTML:} estrutura/semântica (tags semânticas: \texttt{header},
\texttt{nav}, \texttt{main}, \texttt{article}, \texttt{section},
\texttt{footer}).

\textbf{CSS:} apresentação. \textbf{Box model:} content $\to$
\textbf{padding} (interno, dentro da borda) $\to$ \textbf{border} $\to$
\textbf{margin} (externo, fora da borda).

\textbf{Responsividade:} \textbf{media queries} (aplicam estilo conforme a
característica do dispositivo, tipicamente a largura da tela), unidades
relativas (\%, em, rem, vw/vh), layouts flexíveis.

\textbf{Flexbox $\times$ Grid:} o \textbf{Flexbox} é
\textbf{unidimensional} --- distribui itens ao longo de \textbf{um} eixo por
vez (linha \emph{ou} coluna). O \textbf{Grid} é \textbf{bidimensional} ---
define linhas \textbf{e} colunas ao mesmo tempo. Um não substitui o outro:
Grid para a estrutura da página, Flexbox para alinhar o conteúdo dentro de
cada área.

\textbf{Especificidade dos seletores} (quem vence quando duas regras
declaram a mesma propriedade): \texttt{id} (100) $>$ classe / atributo /
pseudoclasse (10) $>$ elemento (1). A ordem de declaração no arquivo só
desempata \textbf{especificidades iguais} --- não vence uma especificidade
maior. (\texttt{!important} passa por cima de tudo, e por isso é
desaconselhado.)

\textbf{\texttt{@import} $\times$ \texttt{<link>}:} os dois trazem uma folha
de estilo externa, mas o \texttt{@import} é resolvido \textbf{dentro do CSS}
e em série (bloqueia o carregamento em cascata), enquanto o \texttt{<link>}
no HTML permite download em paralelo --- por isso \texttt{<link>} é a prática
recomendada. \texttt{@import} \textbf{não} é mecanismo de responsividade (é
distrator recorrente em questão de media query, embora aceite condição de
mídia).

\begin{pegadinha}
\textbf{padding $\times$ margin} (interno $\times$ externo) é o par que a
FGV mais inverte no bloco. Em CSS, some as regras com cuidado:
\texttt{flex-direction: row-reverse} \textbf{inverte} a ordem visual;
\texttt{::before} insere \textbf{antes}, \texttt{::after} \textbf{depois};
\texttt{content: attr(x)} imprime o \textbf{valor} do atributo \texttt{x}.
Resolva no papel, elemento por elemento: aplique o pseudo-elemento, resolva o
\texttt{attr()}, e só então aplique a direção do flex --- a pegadinha está
exatamente na soma dos três. Outras duas inversões do bloco: dizer que
``a \textbf{última} regra declarada sempre vence'' (só vence entre
especificidades \textbf{iguais} --- o seletor de \texttt{id} ganha do de
classe mesmo declarado antes) e trocar \textbf{Flexbox} por \textbf{Grid}
(uni $\times$ bidimensional).
\end{pegadinha}

\section{SPA $\times$ PWA (o coração deste bloco)}

\begin{center}
\begin{tabular}{@{}p{2.6cm}p{4.6cm}p{4.6cm}@{}}
\toprule
& \textbf{SPA (Single Page Application)} & \textbf{PWA (Progressive Web App)} \\
\midrule
O que é & app numa única página; troca conteúdo \textbf{sem recarregar} & site que se comporta como app nativo \\
Chave & roteamento no cliente, sem full reload & \textbf{Service Worker}, offline, \textbf{instalável} \\
Depende de & JS no navegador & recursos web progressivos \\
\bottomrule
\end{tabular}
\end{center}

\textbf{SPA:} uma página, JavaScript atualiza o conteúdo sem recarregar
(React, Angular, Vue).

\textbf{PWA:} usa \textbf{Service Worker} para \textbf{cache, offline e
notificações}, e pode ser \textbf{instalada} no dispositivo como app nativo.

O \textbf{Service Worker} é um script que roda \textbf{em segundo plano},
separado da página, como \textbf{proxy} entre a aplicação e a rede --- é daí
que vêm o cache e o funcionamento offline. Não tem acesso direto ao DOM e
\textbf{exige HTTPS} (só \texttt{localhost} é exceção), justamente porque
interceptar requisições em conexão não cifrada seria um vetor de ataque. O
\textbf{manifest} (\texttt{manifest.json}) é o outro pilar: nome, ícones e
modo de exibição que permitem a \textbf{instalação}.

\begin{pegadinha}
\textbf{Service Worker é da PWA}, não requisito de SPA --- a alternativa que
dá Service Worker à SPA é falsa (caiu exatamente assim na Dataprev 2024).
``SPA instala no SO como nativo'' é característica de \textbf{PWA}. Nenhuma
das duas ``depende exclusivamente'' de framework --- dá para fazer com JS
puro.
\end{pegadinha}

\section{Frameworks}

\begin{center}
\begin{tabular}{@{}p{2.2cm}p{3.4cm}p{3.4cm}p{2.6cm}@{}}
\toprule
& \textbf{React} & \textbf{Angular} & \textbf{Vue} \\
\midrule
Tipo & biblioteca de UI (Meta) & framework completo (Google) & framework progressivo \\
Linguagem & JS/JSX & TypeScript & JS \\
\bottomrule
\end{tabular}
\end{center}

\begin{pegadinha}
A FGV troca React$\leftrightarrow$Angular (biblioteca $\times$ framework
completo) e atribui TypeScript ao React. Em questão de \textbf{React},
cuidado com nomes inventados/parecidos (\texttt{preRender},
\texttt{preloadModule}) no lugar de \texttt{createPortal} (a função certa
para renderizar fora da hierarquia normal do DOM).
\end{pegadinha}

\section{Ajax e comunicação}

\textbf{Ajax:} requisições \textbf{assíncronas} ao servidor sem recarregar a
página (base do comportamento SPA); hoje via \texttt{fetch}/XHR, dados em
JSON.

\section{UX, acessibilidade e usabilidade}

\begin{itemize}
\item \textbf{UX (experiência do usuário)} $\times$ \textbf{UI
(interface)}: UX é a experiência toda; UI é a camada visual.
\item \textbf{Usabilidade} (Nielsen): eficiência, eficácia, satisfação,
facilidade de aprendizado, prevenção de erros.
\item \textbf{Acessibilidade:} \textbf{WCAG} (W3C) --- quatro princípios
\textbf{POUR}: perceptível, operável, compreensível, robusto. Cada critério
de sucesso tem um \textbf{nível de conformidade}: \textbf{A} (mínimo),
\textbf{AA} (o exigido na maioria das normas e contratos públicos) e
\textbf{AAA} (máximo, nem sempre alcançável no site inteiro). No Brasil,
\textbf{eMAG} para governo.
\item \textbf{ARIA} (\textit{Accessible Rich Internet Applications}):
conjunto de atributos --- \texttt{role}, \texttt{aria-label},
\texttt{aria-hidden}, \texttt{aria-live} --- que acrescenta \textbf{semântica}
a elementos \textbf{dinâmicos} ou construídos sem tag nativa, para que
leitores de tela anunciem função, nome e estado. \textbf{Regra número um do
ARIA: não usar ARIA.} Se existe elemento HTML nativo com a semântica certa
(\texttt{<button>}, \texttt{<nav>}, \texttt{<label>}), use o nativo --- o
ARIA é complemento para o que o HTML não cobre, e ARIA mal aplicado piora a
acessibilidade em vez de melhorar.
\item \textbf{Arquitetura de informação:} organização e navegação do
conteúdo.
\item \textbf{CMS} (sistema de gestão de conteúdo): cria/gerencia conteúdo
sem código (WordPress, portais corporativos, workflow editorial).
\end{itemize}

\begin{comosair}[WCAG na prática]
Princípio \textbf{operável}: preferir rótulos/labels bem associados aos
campos e várias formas de entrada além do teclado; fugir de ``submissão
automática ao focar'' (viola previsibilidade) --- é o tipo de alternativa que
a FGV usa como distrator de acessibilidade.
\end{comosair}

\section{Mobile multiplataforma $\times$ nativo}

\begin{itemize}
\item \textbf{Flutter} = \textbf{Dart}; \textbf{React Native} = JS;
\textbf{Ionic} = web; \textbf{Xamarin} = C\#.
\item \textbf{Mobile nativo:} Android = \textbf{Kotlin} (antes Java); iOS =
\textbf{Swift} (antes Objective-C). A FGV troca os pares.
\end{itemize}

\section{Leitura ativa de código --- CSS e React}
\label{sec:leitura-ativa-frontend}

Metade das questões de frontend não pede definição: pede para você
\textbf{prever a renderização} de um trecho de HTML/CSS/JS. Aplique o mesmo
protocolo de leitura ativa do Capítulo \ref{cap:programacao} --- resolva no
papel, propriedade por propriedade, antes de olhar as alternativas.

\subsection{Pegadinhas visuais em CSS}

\begin{conceito}[Some as regras nesta ordem]
\begin{enumerate}
\item \textbf{Box model primeiro.} O tamanho final visível de um elemento é
\texttt{content + padding + border} (\texttt{+ margin} só conta para o espaço
\textbf{ao redor}, não para o tamanho do próprio elemento). Se o enunciado
usa \texttt{box-sizing: border-box}, \texttt{width} já \textbf{inclui}
padding e border --- se usa o padrão (\texttt{content-box}), \texttt{width}
é só o conteúdo e padding/border \textbf{somam por fora}. A FGV monta a
pegadinha exatamente nessa diferença de \texttt{box-sizing}.
\item \textbf{Pseudo-elementos.} \texttt{::before} insere conteúdo
\textbf{antes} do conteúdo real do elemento (mas ainda \textbf{dentro} dele);
\texttt{::after}, \textbf{depois}. Sem \texttt{content: ...} declarado, o
pseudo-elemento \textbf{não aparece}, mesmo com outras propriedades
definidas.
\item \textbf{\texttt{content: attr(x)}.} Imprime o \textbf{valor} do
atributo \texttt{x} do elemento HTML como texto --- se o elemento não tiver
esse atributo, o pseudo-elemento aparece \textbf{vazio}, não gera erro.
\item \textbf{Flex/Grid por último.} \texttt{flex-direction: row-reverse}
inverte a \textbf{ordem visual} dos itens sem alterar a ordem no DOM/HTML;
\texttt{justify-content} e \texttt{align-items} agem nos eixos principal e
transversal --- \textbf{trocar os dois} é o distrator mais comum de layout.
\end{enumerate}
\end{conceito}

\textbf{Exemplo resolvido:} um \texttt{<span data-nivel="urgente">} com CSS
\texttt{span::before \{ content: attr(data-nivel) ": "; font-weight: bold;
\}} $\to$ antes do texto do \texttt{span}, aparece em negrito o literal
\textbf{``urgente: ''} --- se a alternativa disser que aparece o texto
\texttt{"data-nivel"} (o nome do atributo, não o valor), é o distrator
clássico de \texttt{attr()}.

\begin{pegadinha}
A FGV soma 2--3 regras na mesma questão (ex.: \texttt{box-sizing} + pseudo-
elemento + \texttt{flex-direction}) e o erro está na \textbf{combinação}, não
numa regra isolada. Resolva cada camada separadamente antes de somar o
resultado --- não tente ``visualizar direto''.
\end{pegadinha}

\subsection{Escopo e closures em JavaScript}

\begin{conceito}[O item clássico do laço com \texttt{var}]
\texttt{var} tem escopo de \textbf{função} (não de bloco) e sofre
\textit{hoisting}: é içada para o topo da função. Num laço, isso significa
que existe \textbf{uma única} variável, compartilhada por todas as funções
criadas dentro dele --- quando elas forem executadas, todas leem o valor
\textbf{final}.

\begin{center}
\begin{tabular}{@{}p{5.4cm}p{5.4cm}@{}}
\toprule
\textbf{Código} & \textbf{Saída} \\
\midrule
\texttt{for (var i = 0; i < 3; i++)} \newline \texttt{\ \ fs.push(() => console.log(i));} & \textbf{3, 3, 3} --- um \texttt{i} só, que vale 3 ao fim do laço \\
\texttt{for (let i = 0; i < 3; i++)} \newline \texttt{\ \ fs.push(() => console.log(i));} & \textbf{0, 1, 2} --- \texttt{let} cria uma variável \textbf{por iteração} \\
\bottomrule
\end{tabular}
\end{center}
\end{conceito}

\texttt{let} e \texttt{const} (ES6) têm escopo de \textbf{bloco};
\texttt{const} impede a \textbf{reatribuição} da variável, mas \textbf{não}
congela o objeto --- \texttt{const obj = \{\}; obj.x = 1;} é válido.
E \texttt{==} faz \textbf{coerção de tipo} (\texttt{"1" == 1} é
\texttt{true}); \texttt{===} compara valor \textbf{e} tipo.

\begin{pegadinha}
A alternativa que diz que o laço com \texttt{var} imprime \textbf{0, 1, 2} é
o distrator número um --- ela descreve o comportamento do \texttt{let}. O
inverso também aparece: atribuir a \texttt{let} o vazamento de escopo que é
do \texttt{var}. E não confunda com Java, onde o escopo é \textbf{sempre} de
bloco (Capítulo \ref{cap:java}).
\end{pegadinha}

\subsection{Renderização no React}

\begin{itemize}
\item \textbf{\texttt{key} em listas:} React usa a \texttt{key} para
identificar qual item de uma lista mudou, foi adicionado ou removido entre
renderizações. Usar o \textbf{índice do array} como \texttt{key} em lista
que pode reordenar/inserir/remover item \textbf{quebra} esse rastreamento
(o React pode reaproveitar o componente errado para o item errado) --- a
prática recomendada é usar um \textbf{id estável} dos dados, nunca o índice.
\item \textbf{\texttt{useState} é assíncrono e em lote (batched):} chamar o
\textit{setter} de estado não atualiza a variável \textbf{na mesma linha} ---
o valor só reflete a mudança na \textbf{próxima renderização}. Um trecho que
faz \texttt{setContador(contador + 1)} duas vezes seguidas \textbf{não}
soma 2 automaticamente, porque as duas chamadas leem o mesmo \texttt{contador}
``antigo'' dentro do mesmo evento --- para acumular corretamente usa-se a
forma funcional \texttt{setContador(c => c + 1)}.
\item \textbf{Reconciliação (Virtual DOM):} React compara a árvore virtual
nova com a anterior e aplica só as diferenças no DOM real --- não recria a
página inteira a cada render. Criar uma \textbf{função anônima nova a cada
renderização} (ex.: \texttt{onClick={() => ...}} sem memoização) não quebra a
reconciliação, mas pode gerar \textbf{re-renders desnecessários} em
componentes filhos otimizados com \texttt{memo}.
\end{itemize}

\begin{pegadinha}
Alternativa que diz que ``usar o índice do array como \texttt{key} nunca traz
problema'' --- falso quando a lista pode reordenar ou ter itens
inseridos/removidos no meio. Alternativa que diz que \texttt{setState}/
\texttt{useState} atualiza \textbf{imediatamente e de forma síncrona} ---
falso: é assíncrono e em lote dentro do mesmo handler de evento. Confundir
\textbf{Virtual DOM} (estrutura em memória que o React usa para calcular a
diferença) com o \textbf{DOM real} do navegador é outro distrator recorrente.
\end{pegadinha}

\begin{jacaiu}
\textbf{Em prova real da FGV:} leitura de HTML+CSS prevendo a renderização
(pseudo-elementos, \texttt{content}, \texttt{attr()},
\texttt{flex-direction}), a função certa do React
(\texttt{createPortal}, para renderizar fora da hierarquia normal do DOM) e
WCAG na prática (associar rótulos aos campos, princípio \textbf{operável})
--- \textbf{MPU}; SPA $\times$ PWA, com Service Worker/offline/instalável ---
\textbf{Dataprev 2024}.

\textbf{No nosso banco} (previsto pelo edital, ainda não visto na amostra de
provas): box model (padding $\times$ margin, \texttt{box-sizing});
especificidade de seletores; \texttt{Flexbox} $\times$ \texttt{Grid}; media
queries; escopo de \texttt{var} $\times$ \texttt{let} em laço; frameworks
(React $\times$ Angular); \texttt{key} em listas e \texttt{useState}
assíncrono/\textit{batched}; ARIA; UX $\times$ UI; Ajax; CMS.

Rode \texttt{./quiz.py frontend} e \texttt{./quiz.py leitura-codigo}.
\end{jacaiu}

\section{Alta probabilidade / pesquisa extra}

\begin{itemize}
\item \textbf{HTTPS, SSL/TLS} estão no edital do frontend também --- ver
Capítulo \ref{cap:seguranca}.
\item \textbf{SSR $\times$ CSR} (renderização no servidor $\times$ no
cliente); \textbf{hydration}.
\item \textbf{Web Components}, \textbf{micro-frontends} (tendência
arquitetural).
\item \textbf{Core Web Vitals} (performance percebida) e acessibilidade
como requisito legal em serviços públicos (Lei Brasileira de Inclusão).
\end{itemize}

\begin{comosair}
Decore a tabela SPA $\times$ PWA: Service Worker, offline, instalável e
``parece nativo'' são atributos da PWA. Fixe pares mobile: Dart$\to$Flutter,
Kotlin$\to$Android, Swift$\to$iOS.
\end{comosair}
